\subsection{Forcing the bounds at an owner fold}
\label{app:core-owner-folds}

At a fold of one owner, each row can absorb only to that owner's atom.
This determines which bounds all the children must meet and allows
several positive requests to be tested on the same pool.
Fix an owner $r$, a finite set of rows $Q$, upper bounds $A_i\subseteq I$,
and a pool $C$ of permissible child columns in $\Pow^Q$. A finite list
of observations can require an atom $\{a\}$, forbid that atom, or require
the winner $a$ at a specified row $i$. The table gives the resulting
tests on children of a nonempty $\star_r$-fold. An empty entry imposes
no condition.
\[
\begin{array}{ccll}
\text{Observation at }i & \text{Case}& \text{Every child }c & \text{Some child }c \\
\hline
\text{Require }\{a\} & r=a & & c(i)=\{a\} \\
 & r\ne a & c(i)=\{a\} & \\
\text{Forbid }\{a\} & r=a & c(i)\ne\{a\} & \\
 & r\ne a & & c(i)\ne\{a\} \\
\text{Require }a & r=a & & a\in c(i) \\
 & r\ne a & c(i)\ne\{r\} & a\in c(i)
\end{array}
\]
In fact the fold gives $\{r\}$ if a child does and otherwise gives the
union. In particular a foreign singleton requires unanimous children.

Let $B\subseteq C$ be the child columns satisfying every universal
test in the table. For $S\subseteq Q$ put
\[
 C(S)=\{c\in B:\forall i\in S,\ c(i)\subseteq A_i\},\qquad
 F(S)=\{i\in Q:\ r\notin A_i\text{ or no }c\in C(S)
                              \text{ has }c(i)=\{r\}\}.
\]
Starting at $S=\varnothing$, iterate $F$ to its least fixed point $S_*$.
As $S$ grows the pool $C(S)$ shrinks, so $F$ is monotone and the
iteration makes at most $|Q|$ strict increases.

\begin{theorem}[Greatest pool for an owner fold]\label{core:thm:owner-fold-closure}
There is a nonempty finite fold of columns from $C$ meeting all bounds
and observations if and only if $C(S_*)$ is nonempty and each
existential test in the table has a witness in $C(S_*)$. Every such
fold uses only columns in $C(S_*)$. Moreover, any finite set of columns
in $C(S_*)$ can be included in a successful fold when the tests pass.
\end{theorem}
\begin{proof}
Every child of a successful fold lies in $B$. If they all lie in
$C(S)$, then on a row of $F(S)$ there is no child that can absorb to an
allowed atom. The bound at that row therefore requires every child to
obey it. Induction gives containment in $C(S_*)$ and proves necessity.

For sufficiency start with any finite set of columns in $C(S_*)$ and
one default column there. Include one witness per existential test and,
on each row outside $S_*$, one column giving the allowed atom $\{r\}$.
The latter witnesses exist because $F(S_*)=S_*$. At rows in $S_*$
all chosen columns obey the bounds, and outside $S_*$ the chosen atoms
absorb. The universal tests hold on every chosen column, so together
with the witnesses they give the required observations.
\end{proof}

A schedule is a finite word of owners. At its first slot use a
nonempty fold of that owner on any columns supplied by the suffix;
at the empty schedule use a projection. A unary fold can skip a slot.
To decide observations at a slot, compute $S_*$ by one suffix test
per repair query, with the universal observations and the current
bounds in $S$, and with the atom $\{r\}$ required at the queried row.
At stabilization the same suffix tests check nonemptiness and the
individual existential requests, each on the fixed bounds $S_*$.
For a fixed schedule length the number of queries is polynomial in
the explicit observation list. In each slot there are at most
$|Q|(|Q|+1)$ repair queries and one final query per existential request
in addition to the nonemptiness test. All requests in the recursion
are still bounds, atom requirements or prohibitions, or winner
requirements on the same rows.

\begin{corollary}[A uniform schedule for finite samples]\label{core:cor:owner-fold-schedule}
On a nonempty finite set $Q$ of sample rows, repeat any fixed listing
of the $K$ owners $K|Q|$ times. A core term satisfies the sample bounds
and observations exactly when the test for this schedule succeeds.
In particular, on the probes of \cref{core:cor:guard-peeling}, one can
require all positive source bits in the same test and use the same
negative bounds. Along with canonicity this recognizes the prime
source-guard normal forms.
\end{corollary}
\begin{proof}
A witness has some exact tuple of outputs $T_i$ on $Q$. By
\cref{core:thm:subpower-complexity}, the same tuple has a binary witness
in the owner gates of depth at most $\sum_i|T_i|\leq K|Q|$.
For each tree level, use its corresponding sweep of the owners,
placing each gate in the slot for its owner and padding with unary
folds. Conversely, expand any successful finite fold into binary
gates. The guard specialization follows from the probe characterization
in \cref{core:cor:guard-peeling}.
\end{proof}

\paragraph{Changing the measure of depth.}
Three owner slots suffice for an arbitrary nonempty union. For distinct
owners $a,b,c$ and input menus $(X_i)_{i\in J}$, let
$C=\mathop{\star_c}\limits_{j\in J} X_j$ be their nonempty fold at $c$.
Then
\begin{equation}\label{core:eq:three-fold-union}
 \mathop{\star_a}\limits_{i\in J}\bigl(X_i\star_b C\bigr)
      =\bigcup_{i\in J}X_i .
\end{equation}
In fact $X_i\star_b C$ contains $X_i$: erasing one of its winners would
require $C=\{b\}$, but this foreign singleton requires every $X_j$ to
equal $\{b\}$. An operand of the outer fold can equal $\{a\}$ only if
$C=\{a\}$, which requires all inputs to equal $\{a\}$. In that case
the fold is already the input union. Otherwise the outer fold takes the
union, and all its operands are contained in the input union by
provenance. Use a unary $c$-fold for the branch $X_i$ so that all
branches fit the same three slots.

Three slots are necessary for the union of two or more inputs. In two
slots with the same owner the folds flatten, giving that owner's
disjunctive singleton test. For slots $a,b$ with $a\ne b$, the outer
$a$-profile is the disjunction of the unanimity tests on the leaf sets
of its children. For it to be unanimity on all inputs, every child
must use every input. But then the outer $b$-profile is the conjunction
of identical disjunctions over all inputs, so it fails unanimity.

The ordered construction in \cref{core:thm:profiles} already gives the
greatest profile completion. If its output is ambiguous, follow the
option consisting of the whole input set at each slot. In the suffix
after any such options, the current owner's singleton test is still
exactly its prescribed profile. It can choose its winning atomic
coalition when there is one. Conversely, when all other owners choose
the full input set, a successful choice must be contained in the atomic
coalition, making that coalition winning. Thus no slot absorbs along
this branch, so all winners of the final full input union survive.
Together with the singleton tests and provenance, this gives the
completion at every valuation. By \eqref{core:eq:three-fold-union} the
leaf unions require only three additional owner slots.

\begin{theorem}[Compiling completion depth into an owner schedule]
\label{core:thm:completion-fold-schedule}
Fix a listing $\pi$ of the $K\ge3$ owners, and a list $u=(a,b,c)$ of
three distinct owners. Every operation of completion depth at most $d$,
with free unions as in \cref{core:thm:layer-interpolation}, fits the
schedule $u(\pi u u)^d$, of length $(K+6)d+3$. Conversely every word
fitting an $s$-slot schedule has completion depth at most $s$.
On $q\ge1$ sample rows the first conversion can preserve all exact
outputs using a binary word of depth at most
\[
 ((K+6)d+3)\left\lceil\log_2\bigl((3K+1)q+1\bigr)\right\rceil .
\]
\end{theorem}
\begin{proof}
Flatten the free unions above the first completion gates into one
nonempty union. The slots $u$ compute it by
\eqref{core:eq:three-fold-union}. Each first completion gate uses
the listing $\pi$ for the ordered coalition construction and then
$u$ for the leaf unions. Its input expressions have one fewer layer,
so use the remaining schedule there. Projections can skip slots by
unary folds. Conversely, a nonempty owned fold is the greatest
completion of its own profiles, since on ambiguity it takes the union
of its children. Each slot therefore contributes at most one layer.

For the binary bound, require the exact outputs as upper bounds and
winner requests on the $q$ rows. At each slot, the child conditions in
\cref{core:thm:owner-fold-closure} use these same rows, with at most
$3Kq$ distinct observations besides the bounds. To make a successful
fold, keep at most one child for each repairable row, one per existential
observation, and a default child. Apply the same construction to the
suffix requirements on each kept child. Balancing each fold using
associativity gives the displayed bound.
\end{proof}

Taking all valuations on $n$ inputs gives a primitive binary depth of
$O_K(n(d+1))$. Linear dependence on $n$ is necessary already for some
single completion gates. To see this, put unanimity at all but one
owner and let that owner's profile be generated by an arbitrary nonempty
family of middle-sized coalitions. Every such tuple is admissible, and
there are $2^{\binom n{\lfloor n/2\rfloor}}-1$ distinct choices for
$n\ge2$. A primitive word of depth at most $h$ has at most $2^{h+1}-1$
nodes, so there are at most $(n+K+1)^{2^{h+1}}$ such words. Since the
middle binomial coefficient is at least $2^n/(n+1)$, some single gates
need binary depth $n-O_K(\log n)$.

At fixed $K,d$, the schedule test gives a polynomial-time distinction
on explicitly listed samples between feasibility at depth at most $d$
and infeasibility at depth $(K+6)d+3$. For complete prime lists, expand
the probes and use the same test with their winner requests and bounds.
If the prime lists have length $N$, there are at most $2^{O_K(N)}$ rows,
and a schedule of length $s$ takes $2^{O_K(N(s+1))}$ time by the
suffix query bound. The examples of
\cref{core:cor:exact-prime-exponential-depth} require $s$ to grow
faster than any polynomial in $N$, already at bounded prime width.
